\begin{afterword}
\summaryhead

The post argues that very strong evidence is ordinary. When the author tells someone his name, their confidence that it is really his name jumps from about one in a million to near certainty. That means his statement multiplied the odds by roughly twenty million. Wikipedia entries, a friend's claim to juggle, and a computer's output are similar.

From this the author draws two claims. First, you may need less evidence than you think to believe that you are an unusually good trader, or unusually well calibrated. Second, working through Bayes's rule by hand may push people toward estimates that are too timid. The last line turns a familiar maxim around: extraordinary claims require extraordinary evidence, but such evidence may be common.

\responsehead

The opening example is excellent. In four sentences it shows that everyday testimony carries enormous weight, a real and underappreciated point. The post is short and makes one argument, and both are to its credit.

Its weakness is that it never asks why the evidence is strong. The name example works because the claim is narrow and people seldom lie about their names. The trading and calibration examples are different. The claims are broad and flattering, and the evidence (a good year of returns, a good calibration score) is noisy. Here the arithmetic slips as well: 99 to 1, not 200 to 1 (see the margin). The claim that such evidence ``might be easy to get'' comes with no example. A teacher would ask for the second half of the argument to be supported as carefully as the first.

The maxim the post plays on was popularized by Carl Sagan in the late 1970s. The sociologist Marcello Truzzi had written ``extraordinary claims require extraordinary proof'' in 1975, and the idea goes back to David Hume and Pierre-Simon Laplace.

The post's unit, the bit of evidence, has a history of its own. At Bletchley Park in 1940, Alan Turing and I. J. Good measured evidence on a logarithmic scale. They called the unit a ban, and its tenth a deciban. They added decibans from many small clues to decide the daily settings of the German naval Enigma. Turing's codebreakers faced the post's question every day. They had to tell evidence that is strong because it is specific from evidence that only feels strong. The post would have been better for drawing that line.
\end{afterword}
